% !TEX root = ../../main.tex
% Table: FDTD vs FEM vs RCWA comparison for Chapter 14
\begin{table}[htbp]
  \centering
  \small
  \caption{FDTD、FEM 与 RCWA 在 PCSEL 分析中的典型分工。}
  \label{tab:ch14_method_comparison}
  \begin{tabularx}{\textwidth}{@{}p{2.2cm}YYY@{}}
    \toprule
    维度 & FDTD & FEM & RCWA \\
    \midrule
    基本思路 &
    时域推进，直接求解 Maxwell 方程的瞬态响应 &
    由频域弱形式离散，适合复杂边界和本征值问题 &
    平面内 Fourier 展开并沿纵向分层传播，适合周期结构 \\

    适用问题 &
    有限尺寸阵列、宽带响应、近远场提取 &
    冷腔本征模、复频率、局域场与多物理场耦合 &
    能带初筛、反射/透射谱、周期单胞参数扫描 \\

    对周期结构的处理 &
    可以处理，但计算域通常较大 &
    可以处理，但需显式设置 Bloch 或开放边界 &
    最自然，通常作为周期单胞筛选的首选方法 \\

    对复杂几何的处理 &
    正交网格下容易引入 staircase 误差 &
    非结构网格适应性较好，几何保真度高 &
    可处理分层剖面，但对分层近似和周期假设敏感 \\

    计算代价 &
    内存和时间开销通常最高，三维大尺寸问题尤其明显 &
    单次求解较重，但目标频点或目标模分析效率较高 &
    单次扫描较快，适合较大范围参数搜索 \\

    主要输出 &
    时域衰减、宽带响应、近场和远场 &
    本征频率、品质因子、模场分布、局域增益重叠 &
    反射/透射、衍射级次、耦合强度趋势 \\

    在 PCSEL 工作流中的角色 &
    验证有限孔阵、边缘泄漏和实际出光表现 &
    比较候选模、阈值相关损耗与多物理场耦合机制 &
    早期筛选晶格常数、填充因子和扰动参数 \\

    局限性 &
    网格和时间步长受限，长时间衰减问题成本较高 &
    模型搭建和网格质量要求较高 &
    不适合直接描述有限尺寸器件与强非周期扰动 \\
    \bottomrule
  \end{tabularx}

  \vspace{0.6em}
  \parbox{0.97\linewidth}{
    \footnotesize
    对 PCSEL，常见组合流程是：先用 PWEM / RCWA 做带结构和单胞参数初筛，再用 FEM 提取候选模的复频率与场分布，最后用 FDTD 检查有限尺寸、边缘泄漏和远场表现。
  }
\end{table}
